%% kapa: exam=1 round=25 year=2014 period=2 form=A subject=appraisal_law subject_ko=부동산관계법규 seq=32 number=32
%% source: Q-Net 감정평가사 기출문제 (cst00309) artlSeq=5101207 file=extracted/1008/5101207/2014년 제25회 감정평가사 1차시험 기출문제/2교시(A형).hwp
%% conversion: kordoc 4.24.51 -> tools/gen1.py (auto); review=proofread-sonnet; flags=none
\begin{question}{32}
건축법령상 용어에 관한 설명으로 옳은 것은?
\begin{choices}
  \item ‘이전’은 건축물의 주요구조부를 해체하여 같은 대지의 다른 위치로 옮기는 것을 말한다.
  \item 건축물의 피난계단을 증설하는 것은 ‘증축’에 해당한다.
  \item ‘개축’은 건축물이 천재지변으로 멸실된 경우 그 대지에 종전과 같은 규모의 범위에서 다시 축조하는 것을 말한다.
  \item 건축물의 바닥이 지표면 아래에 있는 층으로서 바닥에서 지표면까지 평균높이가 해당 층 높이의 3분의 1인 것은 ‘지하층’에 해당한다.
  \item 바닥(최하층 바닥은 제외)은 ‘주요구조부’에 해당한다.
\end{choices}
\end{question}
\answer{5}
